\documentclass[11pt]{scrartcl}

\usepackage[utf8]{inputenc}
\usepackage[ngerman]{babel}
\usepackage[T1]{fontenc}
\usepackage{graphicx}
\usepackage{amsmath}
\usepackage{float}
\usepackage{caption}
\usepackage{siunitx}
\usepackage{xcolor}
\usepackage{geometry}
\geometry{
    a4paper,
    total={170mm,257mm},
    left=20mm,
    top=20mm,
}
\pagestyle{empty}

\title{\textcolor{red}{Kurzbeschreibung Casino !uv.!}}
\subtitle{Donnerstag Abendprogramm}
\author{}
\date{}
\begin{document}

\maketitle \thispagestyle{empty}
\section{Material}
\begin{itemize}
    \item Drachmen
    \item Roulette
    \item Poker
    \item 4 Gewinnt
    \item Pfeile und Bögen
    \item Juggerwaffen (2 Typgleiche)
    \item Snackboxen
    \item Bingo
    \item Memory
    \item Mocktailzutaten \begin{itemize}
        \item Limetten
        \item Sirup(s)
        \item Eis
    \end{itemize}
\end{itemize}
\section{Grundkonzept}
Die Kinder dürfen frei zwischen den Stationen wechseln und in verschiedenen Spielen ihr Glück auf die Probe stellen. Es gibt zudem Snackboxen und Mocktails für das erspielt Geld zu kaufen (Jeder hat eine begrenzte Anzahl an Snackboxen und Mocktails $\to$ Teamabsprache davor (Vorschlag: 2 Mocktails und 1 Snackbox))
\section{Stationen}
\subsection{Roulette}
\textbf{Einsatz:} 1-5 Drachmen\\
Gespielt wird mit eingeschränkten Setzmöglichkeiten:\begin{itemize}
    \item Gerade/Ungerade (Doppeltes Zurück)
    \item Rot/Schwarz (Doppeltes Zurück)
    \item Hoch/Niedrig (Doppeltes Zurück)
    \item Einzelne Zahl (Fünffaches Zurück)
\end{itemize}
\subsection{Poker}
\textbf{Einsatz:} Freier Einsatz, mindestens 1 Drachme\\
Gespielt wird nach klassischen Poker Regeln. Aussteigen ist nach jedem auszahlen oder nach ausscheiden aus der Runde möglich.
\subsection{4 Gewinnt}
\textbf{Einsatz:} 1 Drachme\\
Gespielt wird klassisches 4 Gewinnt der Gewinner erhält das doppelte von Einsatz.
\subsection{Bogenschießen}
\textbf{Einsatz:} 1-5 Drachmen\\
Es wird 3-mal auf eine Zielscheibe geschossen. Für jeden Treffer $\geq 6$ wird 1 auf den Faktor der Auszahlung addiert:
\begin{itemize}
    \item 0 Treffer: Nichts zurück
    \item 1 Treffer: Einsatz zurück
    \item 2 Treffer: Doppelt zurück
    \item 3 Treffer: Dreifach zurück
\end{itemize}
\subsection{Juggerduell}
\textbf{Einsatz:} 1-5 Drachmen\\
Die Duellanten einigen sich auf einen Einsatz. Der Gewinner erhält den gesamten gemeinsamen Einsatz.
\section{Memory}
\textbf{Einsatz:} 1-5 Drachmen\\Gespielt wird entweder gegen die Bank oder einen weiteren Spieler. Wird Spieler vs. Spieler gespielt, einigen sich die Spieler auf einen Einsatz und der Gewinner erhält den gesammten Einsatz. Wird gegen die Bank gespielt, wird der Einsatz bei einem Sieg verdoppelt. Bei einem Unentschieden gibt es in beiden Fällen den Einsatz zurück.
\section{Bingo}
\textbf{Einsatz:} 3 Drachmen\\Gespielt wird mit mindestens 3 Spielern. Wer als erstes ein Bingo hat gewinnt alle Einsätze, bei mehreren Gewinnern wird der Gewinn geteilt und evtl. auf die nächstkleinere ganze Zahl abgerundet.

\section{Mocktails und Snackboxen}
Es stehen 2-4 Mocktails (genaue Anzahl nach Absprache mit der Küche) zur Auswahl. Weiter gibt es eine Auswahl an einzelnen Snackboxen.\\
Ideen in "\texttt{Hirngespinste.txt}"

\end{document}
